\chapter{\MakeUppercase{Pendahuluan}}
\section{Latar Belakang}
Bagian ini menjelaskan latar belakang yang mendasari dilaksanakannya penelitian. Pada bagian ini, penulis perlu menguraikan alasan, dasar pemikiran, serta faktor-faktor yang mendorong dilakukannya penelitian. Dengan demikian, pembaca dapat memahami mengapa penelitian tersebut perlu dilakukan.

Pada umumnya, penelitian berangkat dari suatu permasalahan yang terjadi pada masa lalu, masa kini, maupun yang diperkirakan akan muncul pada masa mendatang. Oleh karena itu, bagian latar belakang harus menjelaskan permasalahan yang menjadi dasar penelitian serta urgensi untuk mencari solusi atas permasalahan tersebut.

Apabila penelitian didasarkan pada permasalahan yang terjadi di lingkungan sekitar, uraikan kondisi dan permasalahan yang ada secara jelas. Uraian tersebut sebaiknya didukung oleh data hasil survei, laporan resmi, artikel berita, maupun publikasi ilmiah yang relevan.

Apabila penelitian merupakan pengembangan dari suatu sistem atau perangkat yang telah ada, jelaskan kondisi sistem atau perangkat tersebut saat ini beserta keterbatasan atau kekurangan yang menjadi alasan perlunya pengembangan lebih lanjut.

Apabila penelitian merupakan kelanjutan atau pengembangan dari penelitian sebelumnya, jelaskan penelitian-penelitian yang menjadi dasar rujukan. Uraikan tujuan, hasil, serta keterbatasan dari penelitian terdahulu, kemudian jelaskan aspek yang akan dikembangkan, ditingkatkan, atau disempurnakan dalam penelitian yang dilakukan.

Latar belakang penelitian harus memuat hal-hal berikut:

\begin{enumerate}[topsep=0pt,itemsep=0pt,parsep=0pt,leftmargin=*]
\item Gambaran kondisi umum yang relevan dengan topik penelitian.
\item Gambaran kondisi pada bidang atau aspek spesifik yang menjadi fokus penelitian.
\item Permasalahan yang terjadi pada bidang tersebut, baik pada masa lalu, masa kini, maupun potensi permasalahan di masa mendatang.
\item Uraian mengenai permasalahan yang diteliti, meliputi:
\begin{enumerate}[topsep=0pt,itemsep=0pt,parsep=0pt,leftmargin=*]
\item faktor-faktor yang diduga menjadi penyebab permasalahan;
\item karakteristik atau perilaku dari permasalahan yang terjadi;
\item dampak permasalahan terhadap sistem, lingkungan, atau bidang yang lebih luas.
\end{enumerate}
\item Alternatif solusi yang dapat digunakan untuk mengatasi permasalahan tersebut.
\item Solusi yang dipilih beserta alasan pemilihannya.
\item Gambaran singkat mengenai solusi atau pendekatan yang diusulkan dalam penelitian.
\item Pentingnya solusi yang diusulkan, termasuk keunggulan dan dampaknya dibandingkan dengan solusi lain yang telah ada.
\end{enumerate}


\noindent Contoh:

Klasifikasi audio memainkan peran penting dalam sistem pendeteksian anomali kontemporer dengan memberikan {wawasan penting}\label{rev:ru_min_1a} yang sesuai untuk berbagai aplikasi~\cite{Kilickaya2025, LoScudo2023, Barusco2025}. Bidang-bidang ini tidak hanya berlaku untuk lingkungan industri~\cite{Purohit2019} tetapi juga untuk kawasan perkotaan yang dinamis~\cite{Piczak2015, Salamon2014} dan bahkan ekosistem alam yang kompleks, seperti hutan~\cite{Bandara2023}. \dots
\section{Rumusan Masalah}
Bagian ini menjelaskan permasalahan-permasalahan utama yang harus diselesaikan untuk mencapai tujuan penelitian. Setiap rumusan masalah yang diajukan harus memiliki jawaban yang dapat ditemukan pada hasil penelitian, baik dalam bentuk model sistem, analisis, implementasi, pembahasan, lampiran, maupun kesimpulan.

Pada bagian \emph{Rumusan Masalah}, permasalahan yang telah diuraikan pada latar belakang dirumuskan menjadi satu atau beberapa pertanyaan penelitian yang disajikan secara ringkas, jelas, dan terarah. Rumusan masalah harus merepresentasikan inti persoalan yang akan diteliti, bukan kendala teknis atau kesulitan yang mungkin dihadapi peneliti selama proses penelitian, perancangan, implementasi, maupun pengujian sistem.

Contoh rumusan masalah yang baik adalah sebagai berikut:

\begin{enumerate}[topsep=0pt,itemsep=0pt,parsep=0pt,leftmargin=*]
\item Sistem peringatan dini banjir seperti apa yang sesuai untuk diterapkan pada masyarakat di wilayah Dayeuhkolot?
\item Bagaimana merancang dan mengimplementasikan sistem broadcast SMS yang efektif sebagai bagian dari sistem peringatan dini banjir di wilayah Dayeuhkolot?
\end{enumerate}

Sebaliknya, rumusan masalah tidak seharusnya berupa pertanyaan yang berfokus pada detail teknis implementasi, seperti:

\begin{enumerate}[topsep=0pt,itemsep=0pt,parsep=0pt,leftmargin=*]
\item Bagaimana komunikasi antara sensor banjir dan mikrokontroler dilakukan?
\item Bagaimana cara mengintegrasikan modem GPRS dengan mikrokontroler?
\item Bagaimana melakukan konfigurasi modem GPRS untuk mengirimkan SMS secara broadcast?
\end{enumerate}

Rumusan masalah yang telah ditetapkan selanjutnya perlu dijabarkan dalam bentuk uraian yang logis, argumentatif, dan saling terkait. Uraian tersebut harus menggambarkan ruang lingkup penelitian secara jelas, menunjukkan hubungan antarpermasalahan yang diteliti, serta mengarahkan pembahasan secara sistematis menuju tujuan penelitian yang ingin dicapai.


\section{Tujuan dan Manfaat Penelitian}
Bagian ini menjelaskan tujuan yang ingin dicapai melalui penelitian yang dilakukan. Tujuan penelitian harus dirumuskan secara jelas, terukur, dan selaras dengan judul serta rumusan masalah yang telah ditetapkan. Pencapaian tujuan penelitian nantinya harus dapat dibuktikan melalui hasil pembahasan dan dijawab pada bagian kesimpulan.

Selain tujuan, bagian ini juga menjelaskan manfaat yang diharapkan dari hasil penelitian. Manfaat penelitian menggambarkan kegunaan praktis maupun akademis dari solusi, metode, sistem, atau perangkat yang dikembangkan. Hasil penelitian diharapkan dapat memberikan kontribusi dalam menyelesaikan permasalahan yang diteliti, meningkatkan efisiensi dan efektivitas suatu proses, serta menjadi referensi bagi penelitian atau pengembangan selanjutnya.

Hal-hal yang perlu diperhatikan dalam penyusunan tujuan dan manfaat penelitian adalah sebagai berikut:

\begin{itemize}[topsep=0pt,itemsep=0pt,parsep=0pt,leftmargin=*]
\item Tujuan penelitian harus menyatakan secara jelas hal-hal yang ingin dicapai dalam tugas akhir.
\item Tujuan penelitian harus sesuai dan konsisten dengan judul penelitian.
\item Setiap tujuan yang dirumuskan harus dapat dijawab atau dievaluasi pada bagian kesimpulan.
\item Manfaat penelitian harus menjelaskan kegunaan praktis maupun kontribusi yang dihasilkan dari penelitian yang dilakukan.
\end{itemize}


\section{Batasan Masalah}
Bagian ini menjelaskan ruang lingkup penelitian, asumsi yang digunakan, serta kondisi-kondisi yang menjadi batas keberlakuan dari solusi yang diusulkan. Batasan masalah diperlukan untuk memperjelas fokus penelitian sehingga pembahasan dapat dilakukan secara terarah dan sesuai dengan tujuan yang telah ditetapkan.

Batasan masalah juga menjelaskan pada kondisi apa hasil penelitian masih dapat diterapkan dan dianggap valid. Oleh karena itu, batasan yang ditetapkan harus realistis, rasional, serta sesuai dengan kondisi yang sebenarnya. Batasan yang terlalu luas akan menyebabkan ruang lingkup penelitian menjadi tidak terkendali, sedangkan batasan yang terlalu sempit dapat mengurangi relevansi dan manfaat hasil penelitian.

Sebagai contoh, suatu teori atau metode umumnya memiliki batas keberlakuan tertentu, seperti:

\begin{itemize}[topsep=0pt,itemsep=0pt,parsep=0pt,leftmargin=*]
\item Hukum mekanika Newton berlaku pada kondisi kecepatan benda yang jauh lebih rendah dibandingkan kecepatan cahaya.
\item Hukum ekonomi klasik berlaku dengan asumsi *ceteris paribus*, yaitu faktor-faktor lain dianggap tetap.
\end{itemize}

Contoh batasan masalah yang baik adalah sebagai berikut:

\begin{enumerate}[topsep=0pt,itemsep=0pt,parsep=0pt,leftmargin=*]
\item Objek penelitian adalah masyarakat yang berdomisili di Kecamatan Dayeuhkolot.
\item Banjir yang diamati dibatasi pada kejadian yang disebabkan oleh luapan Sungai Citarum.
\item Sistem yang dikembangkan hanya mampu melakukan pengiriman SMS broadcast kepada maksimal 100 nomor tujuan.
\end{enumerate}

Sebaliknya, batasan masalah yang kurang tepat antara lain:

\begin{enumerate}[topsep=0pt,itemsep=0pt,parsep=0pt,leftmargin=*]
\item Beban maksimum mobil listrik yang dikembangkan adalah 10 kg.
\item Sistem hidroponik otomatis hanya dirancang untuk satu tanaman kangkung.
\item Sistem catu daya regeneratif dapat diterapkan pada seluruh jenis kendaraan.
\end{enumerate}

Contoh-contoh tersebut kurang tepat karena batasan yang ditetapkan terlalu sempit, tidak realistis, atau justru terlalu luas sehingga sulit untuk dibuktikan dan dievaluasi dalam penelitian.


\section{Metode Penelitian}
Menjelaskan pendekatan atau metode yang digunakan untuk menyelesaikan pekerjaan dalam \gls{ta}. 

Kegiatan penelitian dilaksanakan melalui beberapa tahapan, yaitu: studi teoritis atau studi literatur untuk membangun landasan konsep, pengukuran empirik untuk memperoleh data lapangan, analisis statistik untuk mengolah dan menginterpretasikan data, simulasi untuk memodelkan dan memprediksi perilaku sistem, perancangan untuk menghasilkan solusi atau prototipe, serta implementasi untuk menerapkan dan menguji hasil rancangan dalam konteks nyata.

\section{Jadwal Pelaksanaan}
Bagian ini memuat rencana pelaksanaan tugas akhir yang disusun dalam bentuk tahapan kegiatan beserta target waktu penyelesaiannya. Jadwal pelaksanaan berfungsi sebagai pedoman dalam mengelola, memantau, dan mengevaluasi kemajuan penelitian selama periode pengerjaan tugas akhir.

Setiap tahapan kegiatan perlu dilengkapi dengan \emph{milestone}, yaitu target capaian yang dapat digunakan sebagai indikator keberhasilan suatu tahap pekerjaan. \emph{Milestone} harus dirumuskan secara jelas, terukur, dan dapat diverifikasi sehingga dapat digunakan sebagai acuan dalam menilai perkembangan penelitian.

Jadwal pelaksanaan yang disusun akan menjadi dasar pemantauan kemajuan tugas akhir, baik oleh mahasiswa maupun dosen pembimbing. Oleh karena itu, jadwal dan \emph{milestone} yang ditetapkan harus realistis, sesuai dengan ruang lingkup penelitian, serta mempertimbangkan ketersediaan sumber daya dan waktu yang dimiliki.

Contoh jadwal pelaksanaan dan \emph{milestone} dapat dilihat pada Tabel~\ref{tab:jadwal}.

\begin{table}[!ht]
    \centering
    \caption{Contoh Jadwal Pelaksanaan dan \textit{Milestone}}
    \label{tab:jadwal}
    \begin{tabular}{|c|p{0.25\textwidth}|c|c|p{0.30\textwidth}|}
        \hline
        No. & Tahapan Kegiatan & Durasi & Target Selesai & \textit{Milestone} \\
        \hline
        1 & Studi literatur dan analisis kebutuhan & 2 minggu & 22 Januari 2027 & Kebutuhan sistem dan spesifikasi penelitian terdefinisi \\
        \hline
        2 & Perancangan sistem & 2 minggu & 5 Februari 2027 & Diagram blok, arsitektur sistem, dan spesifikasi masukan-keluaran selesai \\
        \hline
        3 & Implementasi dan integrasi sistem & 4 minggu & 5 Maret 2027 & Prototipe sistem versi pertama selesai \\
        \hline
        4 & Pengujian dan evaluasi sistem & 2 minggu & 19 Maret 2027 & Hasil pengujian dan analisis performa sistem tersedia \\
        \hline
        5 & Penyusunan laporan tugas akhir & 4 minggu & 16 April 2027 & Draf laporan tugas akhir selesai \\
        \hline
        6 & Revisi dan finalisasi laporan & 2 minggu & 30 April 2027 & Naskah tugas akhir siap diajukan \\
        \hline
    \end{tabular}
\end{table}
